\maketitle

\begin{abstract}
Starting from a $\sigma$-additive probability measure on an observable
algebra~\cite{mahon_paper1}, this paper derives two structures that are often
assumed: the dynamics of prediction, and the recoverability of the state space.

Making prediction precise forces the introduction of the predictive kernel
$\Pi_Q : \alpha \to \mathcal{P}(\beta)$, the regular conditional distribution of
a future observation $F$ given a current observation $Q$.  Composing with $Q$ gives
the minimal predictive state map $Q_* : \Omega \to \mathcal{P}(\beta)$.  Indexing
over time, the kernels compose via Chapman--Kolmogorov — derived, not assumed —
yielding a Markov semigroup $\{K_t\}$ and the Koopman--Perron duality
$\int K_t g\, d\mu = \int g\, d(\mu \star \Pi_t)$.
When the kernels are Dirac measures, $K_t$ collapses to the classical Koopman
operator.

The construction produces a predictive state space $\mathcal{P}(\beta)$.
In a measure-preserving system $(X, \mathcal{B}, \mu, T)$ with
$Q_t = h \circ T^t$, the map $Q_*$ becomes the delay map
$\Phi_h(x) = (h(T^n x))_{n \in \mathbb{Z}}$.
The reconstruction question is then internal to the same framework: whether
this predictive state is faithful — whether the prediction machinery sees all
of $X$, or only a proper quotient.

The answer is a three-way equivalence — the reconstruction theorem.
State-space recovery ($\mathcal{O}_h = \mathcal{B}$ mod $\mu$) holds if and only
if the delay algebra $\mathcal{A}_h$ is $L^2$-dense, if and only if $\Phi_h$ is a
measure-theoretic embedding.  The bridge between the first two conditions is a
general lemma: for any subalgebra $\mathcal{A} \subseteq L^\infty$ containing $1$,
the $L^2$-closure of $\mathcal{A}$ equals $L^2(X, \sigma(\mathcal{A}), \mu)$.
When reconstruction holds, the compact Stone space~\cite{mahon_paper1} —
built as a completion of the sample space for deriving $\sigma$-additivity —
is identified with $X$ itself.
\end{abstract}

\begin{quote}\small
\textbf{2020 Mathematics Subject Classification.}
37A30 (primary); 47D06, 60J05, 37A05, 28A60, 06E15, 46E30 (secondary).

\textbf{Keywords.}
Koopman operator, predictive kernel, minimal predictive state, Markov semigroup,
Koopman--Perron duality, observable dynamical system, reconstruction theorem,
delay embedding, observable algebra, density bridge, measure-theoretic Takens,
Stone duality, measure-preserving system.
\end{quote}

% ====================================================================
\section{Introduction}
\label{II:sec:intro}

The classical approach to dynamics, following Koopman~\cite{Koopman1931},
describes how observable functions evolve under a known transformation $T$: the
operator $U_T f = f \circ T$ acts by precomposing with $T$ as a primitive.  This
paper asks two questions that precede this picture.

The first is: given a current observation $Q : \Omega \to \alpha$ and a future
observation $F : \Omega \to \beta$, what does an observer predict about $F$ given
$Q$?  The dynamics, if there are any, should follow from the answer — not precede
it.  Making prediction precise immediately requires a $\sigma$-additive probability
measure on the observable algebra.  That measure is not an extra assumption: a coherent family of finitely additive
charges on a directed system of observable algebras extends to a genuine probability
measure exactly when it satisfies collective exhaustion~\cite{mahon_paper1}, the
condition ruling out persistent mass on globally vanishing events.

Once the measure exists, the canonical answer to the prediction question is the
regular conditional distribution of $F$ given $Q$: the unique Markov kernel
$\Pi_Q : \alpha \to \mathcal{P}(\beta)$ realising $P(F \in \cdot \mid Q = q)$
simultaneously across all outcomes.  Its existence rests on the Rokhlin
disintegration theorem~\cite{Rokhlin1952}, which requires $\sigma$-additivity.
Composing with $Q$ gives the minimal predictive state map
$Q_* : \Omega \to \mathcal{P}(\beta)$, the minimal sufficient statistic for future
prediction.  Indexing over time, the kernels $\{\Pi_t\}$ compose; their composition
rule is the Chapman--Kolmogorov equation, derived from temporal coherence rather than
assumed.  From it emerge a semigroup of Markov operators $\{K_t\}$ and the
Koopman--Perron duality connecting them with the pushforward operators.  When the
kernels are Dirac measures, the construction collapses to the classical Koopman
picture.

The construction produces a predictive state space $\mathcal{P}(\beta)$ and a
stationary distribution $\nu_{Q_*}$.  In a measure-preserving system
$(X, \mathcal{B}, \mu, T)$ with $Q_t = h \circ T^t$ for some $h \in L^\infty$,
the map $Q_*$ specialises to the delay map
$\Phi_h(x) = (h(T^n x))_{n \in \mathbb{Z}}$.  The reconstruction question is
therefore not a separate problem: it asks whether this predictive state is
faithful — whether the prediction machinery sees all of $X$, or only a proper
quotient.  The answer is a three-way equivalence: the $\sigma$-algebraic
condition $\mathcal{O}_h = \mathcal{B}$ modulo $\mu$, the analytic condition
that the delay algebra $\mathcal{A}_h$ is $L^2$-dense, and the geometric
condition that $\Phi_h$ is a measure-theoretic embedding are all equivalent.

The key to the equivalence is the density bridge lemma: for any subalgebra
$\mathcal{A} \subseteq L^\infty(X,\mu)$ containing the constant $1$, the
$L^2$-closure of $\mathcal{A}$ equals $L^2(X, \sigma(\mathcal{A}), \mu)$.  The
proof goes via the functional monotone class theorem~\cite{Dynkin}.  This is
what lets the $\sigma$-algebraic condition ($\mathcal{O}_h = \mathcal{B}$
mod $\mu$) and the analytic condition ($\mathcal{A}_h$ $L^2$-dense) pass freely
between each other — they are not two separate conditions but two descriptions
of the same constraint.

When reconstruction holds, a compact space introduced in~\cite{mahon_paper1} for
an entirely different purpose — the Stone space $\mathrm{St}(\mathcal{O}_h)$ of
the observable algebra, built as a technical completion of the sample space for
deriving $\sigma$-additivity — turns out to be measure-theoretically isomorphic to
$X$ itself.  The programme that began by seeking a compact ambient space for the
measure ends by identifying it with the state space being observed.

The paper is self-contained modulo the probability measure.  Sections~\ref{II:sec:prediction}
and~\ref{II:sec:dynamics} develop the predictive structure from a probability space
$(\Omega, \mathcal{F}, P)$.  Section~\ref{II:sec:transition} makes the connection explicit: specialising $Q_*$
to a measure-preserving system yields the delay map, and the reconstruction
question becomes whether that map is injective modulo null sets.
Sections~\ref{II:sec:density-bridge} and~\ref{II:sec:stone} carry out the
reconstruction, proving the density bridge and the reconstruction theorem.

% ====================================================================
\section{Prediction and the minimal predictive state}
\label{II:sec:prediction}

Fix a probability space $(\Omega, \mathcal{F}, P)$ and two measurable maps:
\begin{itemize}
  \item $Q : \Omega \to \alpha$, where $(\alpha, \mathcal{A})$ is a measurable space;
  \item $F : \Omega \to \beta$, where $(\beta, \mathcal{B})$ is a standard Borel space.
\end{itemize}
The standard Borel hypothesis on $\beta$ is used once: it is the hypothesis under
which the Rokhlin disintegration theorem~\cite{Rokhlin1952} guarantees the existence
of a regular conditional distribution.  All other results hold in greater generality.

We write $\mathcal{P}(\beta)$ for the space of probability measures on $\beta$,
equipped with its canonical (Giry) $\sigma$-algebra.  The joint pushforward
$\rho := P \circ (Q, F)^{-1}$ on $\alpha \times \beta$ has marginals
$\rho_\alpha = P \circ Q^{-1}$ and $\rho_\beta = P \circ F^{-1}$.

The regular conditional distribution of $F$ given $Q$ is the unique (up to
$P \circ Q^{-1}$-null sets) Markov kernel $\Pi_Q : \alpha \to \mathcal{P}(\beta)$
satisfying
\[
  \int_A \Pi_Q(q, B)\, d(P \circ Q^{-1})(q) \;=\; P(Q \in A,\; F \in B)
\]
for all measurable $A \subseteq \alpha$ and $B \subseteq \beta$.

\begin{definition}[Predictive kernel]
\label{II:def:predictive-kernel}
The \emph{predictive kernel} is $\Pi_Q$, the disintegration kernel of
$\rho = P \circ (Q, F)^{-1}$ over its first marginal:
$\Pi_Q(q, A) = P(F \in A \mid Q = q)$.
\end{definition}

The kernel is covariant under measurable maps: if $\pi : \alpha \to \alpha'$ is
measurable and $Q' = \pi \circ Q$, then $\Pi_Q(q, \cdot) = \Pi_{Q'}(\pi(q), \cdot)$
for $(P \circ Q^{-1})$-almost every $q$, by uniqueness of disintegration.

Composing $\Pi_Q$ with $Q$ gives the map sending each sample point to its
conditional law.

\begin{definition}[Minimal predictive state map]
\label{II:def:minimal-predictive}
The \emph{minimal predictive state map} is
\[
  Q_* : \Omega \to \mathcal{P}(\beta), \qquad Q_*(\omega) \;=\; \Pi_Q(Q(\omega), \cdot).
\]
\end{definition}

Since $\Pi_Q$ is a Markov kernel it is measurable, and $Q_*$ inherits
measurability by composition.

\begin{theorem}[Predictive factorization]
\label{II:thm:predictive-factorization}
For every bounded measurable $g : \beta \to \mathbb{R}$:
\[
  E[g(F) \mid \sigma(Q)] \;=\; \int g\, dQ_*(\cdot)
  \quad P\text{-a.e.}
\]
\end{theorem}
\begin{proof}
Let $h(\omega) := \int g\, dQ_*(\omega)$.  We verify the two defining properties.

\emph{Measurability.}  The map $q \mapsto \int g\, d\Pi_Q(q, \cdot)$ is
strongly measurable; since $Q_* = \Pi_Q(Q(\cdot), \cdot)$, the map $h$ is
$\sigma(Q)$-measurable.

\emph{Set-integral identity.}  For $S = Q^{-1}(A)$:
\begin{align*}
  \int_S h\, dP
  &= \int_A \int g\, d\Pi_Q(q, \cdot)\, d(P \circ Q^{-1})(q)
  \quad\text{(pushforward)} \\
  &= \int_{A \times \beta} g(f)\, d\rho(q, f)
  \quad\text{(disintegration)} \\
  &= \int_S g(F)\, dP.
\end{align*}
The conclusion follows from a.e.\ uniqueness of conditional expectation.
\end{proof}

\begin{corollary}[Predictive sufficiency]
\label{II:cor:sufficiency}
$E[g(F) \mid \sigma(Q)] = E[g(F) \mid \sigma(Q_*)]$ $P$-a.e.
\end{corollary}
\begin{proof}
$\sigma(Q_*) \subseteq \sigma(Q)$, and the conditional expectation given
$\sigma(Q)$ is already $\sigma(Q_*)$-measurable by the theorem.
\end{proof}

The operator $(K_{Q_*} g)(q_*) = \int g\, dq_*$ on $\mathcal{P}(\beta)$
recasts the theorem as: $E[g(F) \mid \sigma(Q)] = (K_{Q_*} g) \circ Q_*$
$P$-a.e.

% ====================================================================
\section{Dynamics and duality}
\label{II:sec:dynamics}

To introduce dynamics, index over a discrete time horizon.  For each
$t \in \mathbb{N}$, let $\Pi_t : \gamma \to \mathcal{P}(\gamma)$ be a Markov
kernel on the state space $\gamma$, with $\Pi_0(q, \cdot) = \delta_q$.
A family $\{\Pi_t\}$ satisfying the Chapman--Kolmogorov equation
\[
  \Pi_{t+s}(q, \cdot) \;=\; \int_\gamma \Pi_t(q', \cdot)\, d\Pi_s(q, q')
  \quad\text{(i.e., } \Pi_{t+s} = \Pi_t \circ_\kappa \Pi_s\text{)}
\]
is a \emph{Markov semigroup kernel}.  The associated \emph{Markov operators} are
\[
  (K_t g)(q) \;:=\; \int g\, d\Pi_t(q, \cdot), \qquad g : \gamma \to \mathbb{R}.
\]

Each $K_t$ is positive, unital, and contractive on bounded measurable functions,
with $K_0 = \mathrm{id}$; positivity, unitality, and contractiveness are immediate
from the Markov property.

\begin{proposition}[Semigroup property]
\label{II:prop:semigroup}
$K_{t+s} g = K_t(K_s g)$ for all bounded measurable $g$.
\end{proposition}
\begin{proof}
$(K_{t+s} g)(q)
= \int g\, d(\Pi_t \circ_\kappa \Pi_s)(q, \cdot)
= \int \bigl(\int g\, d\Pi_t(q', \cdot)\bigr)\, d\Pi_s(q, q')
= \int (K_t g)(q')\, d\Pi_s(q, q')
= (K_t(K_s g))(q)$,
where the exchange of integrals is justified by Bochner--Fubini.
\end{proof}

The pushforward of a measure $\mu$ on $\gamma$ through the kernel is
$P_t^* \mu := \mu \star \Pi_t = \int \Pi_t(q, \cdot)\, d\mu(q)$.

\begin{theorem}[Koopman--Perron duality]
\label{II:thm:koopman-perron}
For every bounded measurable $g$ and finite measure $\mu$:
\[
  \int (K_t g)\, d\mu \;=\; \int g\, d(P_t^* \mu).
\]
\end{theorem}
\begin{proof}
$\int g\, d(P_t^* \mu)
= \int g\, d(\mu \star \Pi_t)
= \int_\gamma \int g\, d\Pi_t(q, \cdot)\, d\mu(q)
= \int (K_t g)\, d\mu$,
by Bochner--Fubini.
\end{proof}

When $\Pi_t(q, \cdot) = \delta_{\varphi_t(q)}$ for a measurable map
$\varphi_t : \gamma \to \gamma$, the stochastic picture collapses.

\begin{proposition}[Deterministic limit]
\label{II:prop:deterministic}
If $\Pi_t(q, \cdot) = \delta_{\varphi_t(q)}$, then $K_t g = g \circ \varphi_t$
and $\varphi_{t+s} = \varphi_t \circ \varphi_s$.
\end{proposition}
\begin{proof}
$(K_t g)(q) = \int g\, d\delta_{\varphi_t(q)} = g(\varphi_t(q))$.
The flow law follows from
$\delta_{\varphi_{t+s}(q)} = (\Pi_t \circ_\kappa \Pi_s)(q, \cdot)
= \delta_{\varphi_t(\varphi_s(q))}$
and injectivity of the Dirac embedding.
\end{proof}

In the deterministic case, $K_t g = g \circ \varphi_t = U_T^t g$: the Markov
operator specialises to the classical Koopman operator.  The
Koopman--Perron duality $\int K_t g\, d\mu = \int g\, d(P_t^* \mu)$ becomes
the change-of-variables formula $\int g \circ T^t\, d\mu = \int g\, d(T^t_* \mu)$.

The full structure produced by the prediction construction is captured in one
definition.

\begin{definition}[Observable dynamical system]
\label{II:def:obs-dyn-sys}
An \emph{observable dynamical system} is a triple
$(\mathcal{P}(\beta),\, \nu_{Q_*},\, \{K_t\}_{t \in \mathbb{N}})$
where $\mathcal{P}(\beta)$ is the predictive state space (the image of $Q_*$),
$\nu_{Q_*} = P \circ Q_*^{-1}$ is the pushforward of $P$ along $Q_*$, and
$\{K_t\}$ is the Markov semigroup on $\mathcal{P}(\beta)$.
\end{definition}

% ====================================================================
\section{From prediction to reconstruction}
\label{II:sec:transition}

Setting $Q_t = h \circ T^t$ in a measure-preserving system
$(X, \mathcal{B}, \mu, T)$ with $h \in L^\infty$ specialises the construction.
Under this choice, $Q_*(\omega) = (h(T^n \omega))_{n \in \mathbb{Z}}$ is the
delay map $\Phi_h(x) = (h(T^n x))_{n \in \mathbb{Z}}$, and the observable
dynamical system is the triple $(\Phi_h(X), \mu \circ \Phi_h^{-1}, \{K_t\})$.

The predictive state space $\Phi_h(X)$ is a subset of $\mathbb{R}^\mathbb{Z}$
encoding the dynamical information that $h$ can see.  Whether it is a faithful
image of $X$ depends on how much the delay orbit
$\{h \circ T^n : n \in \mathbb{Z}\}$ collectively distinguishes points of $X$.
Injectivity of $\Phi_h$ modulo $\mu$-null sets is equivalent to the observable
$\sigma$-algebra
\[
  \mathcal{O}_h = \sigma\!\bigl(\{h \circ T^n : n \in \mathbb{Z}\}\bigr)
\]
equalling $\mathcal{B}(X)$ modulo $\mu$.

The Koopman--Perron duality (Theorem~\ref{II:thm:koopman-perron}) connects this to a
function-space question: the observable $\sigma$-algebra $\mathcal{O}_h$ determines
the Markov semigroup through the algebra it generates, and the reconstruction question
is precisely whether that algebra is large enough to recover all of $\mathcal{B}$.
The next two sections answer this question.

% ====================================================================
\section{Density bridge and reconstruction}
\label{II:sec:density-bridge}

Throughout this section and the next, $(X, \mathcal{B}, \mu)$ is a standard
probability space, $T : X \to X$ is a measurable, $\mu$-preserving invertible
bimeasurable bijection with $T_*\mu = \mu$, and $h \in L^\infty(X, \mu)$.

\begin{definition}[Observable algebra]
\label{II:def:obs-algebra}
The \emph{observable algebra} generated by $h$ is
\[
  \mathcal{O}_h \;:=\; \sigma\!\bigl(\{h \circ T^n : n \in \mathbb{Z}\}\bigr)
  \;\subseteq\; \mathcal{B}(X).
\]
The \emph{delay algebra} $\mathcal{A}_h$ is the smallest subalgebra of
$L^\infty(X, \mu)$ containing $1$ and every $h \circ T^n$; concretely, finite
sums $\sum_\alpha c_\alpha \prod_k (h \circ T^{n_k})^{e_k}$.
\end{definition}

\begin{definition}[Delay map]
\label{II:def:delay-map}
The \emph{delay map} is
\[
  \Phi_h : X \to \mathbb{R}^\mathbb{Z},
  \qquad
  \Phi_h(x) \;=\; \bigl(h(T^n x)\bigr)_{n \in \mathbb{Z}}.
\]
Since $\pi_n \circ \Phi_h = h \circ T^n$ for each coordinate projection $\pi_n$,
the map $\Phi_h$ is measurable and
$\mathcal{O}_h = \Phi_h^{-1}(\mathcal{B}(\mathbb{R}^\mathbb{Z}))$.
Thus $\mathcal{O}_h = \mathcal{B}$ mod $\mu$ is precisely the condition that
$\Phi_h$ is a measure-theoretic embedding.
\end{definition}

The central lemma connects the Boolean algebra level to the $L^2$ level.
It holds for any subalgebra of $L^\infty$.

\begin{lemma}[Density bridge]
\label{II:lem:density-bridge}
Let $(X, \mathcal{B}, \mu)$ be a probability space and let
$\mathcal{A} \subseteq L^\infty(X, \mu)$ be a subalgebra containing the
constant function~$1$.  Then:
\[
  \overline{\mathcal{A}}^{L^2}
  \;=\;
  L^2(X,\, \sigma(\mathcal{A}),\, \mu).
\]
Consequently, $\overline{\mathcal{A}}^{L^2} = L^2(X, \mu)$ if and only if
$\sigma(\mathcal{A}) = \mathcal{B}$ modulo $\mu$.
\end{lemma}

\begin{proof}
Write $\mathcal{M} := \sigma(\mathcal{A})$ and $V := \overline{\mathcal{A}}^{L^2}$.

\textbf{Step 1: $V \subseteq L^2(X, \mathcal{M}, \mu)$.}
Every $f \in \mathcal{A}$ is $\mathcal{A}$-measurable, hence $\mathcal{M}$-measurable.
So $\mathcal{A} \subseteq L^2(X, \mathcal{M}, \mu)$.  Since $L^2(X, \mathcal{M}, \mu)$
is a closed subspace of $L^2(X, \mathcal{B}, \mu)$, taking the closure gives
$V \subseteq L^2(X, \mathcal{M}, \mu)$.

\textbf{Step 2: $L^2(X, \mathcal{M}, \mu) \subseteq V$.}
It suffices to show that every indicator $\mathbf{1}_E$ with $E \in \mathcal{M}$
lies in $V$.  By the functional monotone class theorem~\cite{Dynkin}, the class
\[
  \mathcal{H}
  \;:=\;
  \bigl\{f \in L^\infty(X, \mathcal{M}, \mu) : f \in V\bigr\}
\]
is a vector space containing $\mathcal{A}$ and closed under bounded monotone
pointwise limits (since if $f_n \in V$ are uniformly bounded and $f_n \to f$
pointwise a.e.\ with $f \in L^\infty$, then $f_n \to f$ in $L^2$ by dominated
convergence, so $f \in V$).  Since $\mathcal{A}$ is an algebra containing $1$
and $\sigma(\mathcal{A}) = \mathcal{M}$, the monotone class theorem gives
$\mathcal{H} = L^\infty(X, \mathcal{M}, \mu)$.  In particular every
$\mathcal{M}$-indicator $\mathbf{1}_E \in V$.

Since the $\mathcal{M}$-simple functions are dense in $L^2(X, \mathcal{M}, \mu)$
and $V$ is closed, $L^2(X, \mathcal{M}, \mu) \subseteq V$.

\textbf{Conclusion.}
Steps~1 and~2 give $V = L^2(X, \mathcal{M}, \mu)$.  The final equivalence is
immediate: $V = L^2(X, \mathcal{B}, \mu)$ iff $L^2(X, \mathcal{M}, \mu) =
L^2(X, \mathcal{B}, \mu)$ iff $\mathcal{M} = \mathcal{B}$ mod $\mu$.
\end{proof}

\begin{remark}[Role of the algebra structure]
\label{II:rem:algebra-necessary}
The algebra structure of $\mathcal{A}$ is essential.  The corresponding
statement for the closed \emph{linear} span $\mathcal{H}_{\mathcal{A}} =
\overline{\mathrm{span}}(\mathcal{A})$ is false in general:
$\mathcal{H}_{\mathcal{A}} = L^2$ does not imply
$\sigma(\mathcal{A}) = \mathcal{B}$.

For example, take $X = [0,1]$, $\mu = $ Lebesgue, $T(x) = 2x \bmod 1$ (the
doubling map), and $h(x) = x$.  The orbit $\{h \circ T^n\}(x) = \{2^n x \bmod 1\}$
generates $\mathcal{B}([0,1])$ as a $\sigma$-algebra (the dyadic structure separates
points).  But $\mathrm{span}\{2^n x \bmod 1 : n \geq 0\}$ is a proper subspace of
$L^2[0,1]$: it does not contain $x^2$.  So the linear span can be a proper dense
subspace even when $\sigma(\mathcal{A}) = \mathcal{B}$.

In the other direction, $h$ cyclic for $U_T$ (linear span dense) always implies
$\sigma(\mathcal{A}_h) = \mathcal{B}$ mod $\mu$ — see
Corollary~\ref{II:cor:cyclic-implies-reconstruction}.
\end{remark}

\begin{theorem}[Reconstruction theorem]
\label{II:thm:reconstruction}
Let $(X, \mathcal{B}, \mu, T)$ be an invertible measure-preserving system and
$h \in L^\infty(X, \mu)$.  The following are equivalent:
\begin{enumerate}[label=(\roman*)]
  \item $\mathcal{O}_h = \mathcal{B}$ modulo $\mu$;
  \item $\mathcal{A}_h$ is dense in $L^2(X, \mu)$;
  \item $\Phi_h : X \to \mathbb{R}^\mathbb{Z}$ is a measure-theoretic embedding:
    $\Phi_h^{-1}(\mathcal{B}(\mathbb{R}^\mathbb{Z})) = \mathcal{B}$ modulo $\mu$.
\end{enumerate}
\end{theorem}

\begin{proof}
\emph{(i) $\Leftrightarrow$ (ii).}
Apply the density bridge (Lemma~\ref{II:lem:density-bridge}) with
$\mathcal{A} = \mathcal{A}_h$ and $\sigma(\mathcal{A}_h) = \mathcal{O}_h$:
$\overline{\mathcal{A}_h}^{L^2} = L^2(X, \mu)$ iff $\mathcal{O}_h = \mathcal{B}$
mod $\mu$.

\emph{(i) $\Leftrightarrow$ (iii).}
By Definition~\ref{II:def:delay-map},
$\mathcal{O}_h = \Phi_h^{-1}(\mathcal{B}(\mathbb{R}^\mathbb{Z}))$.
So $\mathcal{O}_h = \mathcal{B}$ mod $\mu$ iff
$\Phi_h^{-1}(\mathcal{B}(\mathbb{R}^\mathbb{Z})) = \mathcal{B}$ mod $\mu$,
which is exactly condition~(iii).
\end{proof}

\begin{remark}[What ``measure-theoretic embedding'' means]
\label{II:rem:embedding}
Condition~(iii) says that $\Phi_h$ is \emph{almost everywhere injective} and
the pullback $\sigma$-algebra is all of $\mathcal{B}$.  More precisely: since
$\Phi_h^{-1}(\mathcal{B}(\mathbb{R}^\mathbb{Z})) = \mathcal{B}$ mod $\mu$, for
any two sets $A, A' \in \mathcal{B}$ with $\mu(A \triangle A') > 0$ there exists
a Borel set $E \subseteq \mathbb{R}^\mathbb{Z}$ with
$\mu(\Phi_h^{-1}(E) \triangle A) = 0$ and
$\mu(\Phi_h^{-1}(E) \triangle A') > 0$.  In measure-theoretic language,
$\Phi_h$ induces an isomorphism of measure algebras
$(X/\mu, \mathcal{B}/\mu) \cong (\mathbb{R}^\mathbb{Z}/(\Phi_h)_*\mu,
\mathcal{B}(\mathbb{R}^\mathbb{Z})/(\Phi_h)_*\mu)$.

The $T$-invariance of $\mu$ implies that $(\Phi_h)_*\mu$ is invariant under
the bilateral shift $\sigma$ on $\mathbb{R}^\mathbb{Z}$, since
$(\Phi_h(Tx))_n = h(T^{n+1}x) = (\sigma(\Phi_h(x)))_n$ for each $n$.  The
dynamical system $(X, T, \mu)$ is therefore isomorphic to
$(\mathbb{R}^\mathbb{Z}, \sigma, (\Phi_h)_*\mu)$ restricted to the image of
$\Phi_h$.
\end{remark}

% ====================================================================
\section{Cyclic vectors and Stone space identification}
\label{II:sec:stone}

\begin{definition}[Cyclic vector]
\label{II:def:cyclic}
The function $h \in L^2(X, \mu)$ is a \emph{cyclic vector} for $U_T$ if
\[
  \mathcal{H}_h
  \;:=\;
  \overline{\mathrm{span}}\!\bigl\{U_T^n h : n \in \mathbb{Z}\bigr\}
  \;=\; L^2(X, \mu).
\]
\end{definition}

\begin{corollary}[Cyclic implies reconstruction]
\label{II:cor:cyclic-implies-reconstruction}
If $h \in L^\infty(X, \mu)$ is a cyclic vector for $U_T$, then
$\mathcal{O}_h = \mathcal{B}$ modulo $\mu$.
\end{corollary}

\begin{proof}
Since $\mathcal{A}_h$ contains $\mathrm{span}\{h \circ T^n\} =
\mathrm{span}\{U_T^n h\}$ as a subset, $\overline{\mathcal{A}_h}^{L^2}
\supseteq \mathcal{H}_h = L^2$.  By the density bridge, $\mathcal{O}_h =
\mathcal{B}$ mod $\mu$.
\end{proof}

\begin{remark}[The implication is strict]
\label{II:rem:strict}
The converse of Corollary~\ref{II:cor:cyclic-implies-reconstruction} fails in
general: $\mathcal{O}_h = \mathcal{B}$ does not imply $h$ is cyclic.
The doubling map example of Remark~\ref{II:rem:algebra-necessary} witnesses this.

The cyclic vector condition asks for density of the \emph{linear span} of the
orbit; reconstruction asks only for density of the \emph{algebra}.  For a bounded
$h$, the linear span is contained in the algebra, so cyclicity is the stronger
condition.

The two conditions coincide precisely when $U_T$ has \emph{simple
$L^2$-spectrum}: $L^2(X,\mu)$ is a cyclic $U_T$-module, meaning there exists
some cyclic vector.  In this case, every generating observation $h$ (satisfying
$\mathcal{O}_h = \mathcal{B}$) is automatically cyclic.  Simple spectrum is
a generic property of ergodic transformations in the Baire category sense
\cite{Halmos1944} but is not universal.
\end{remark}

The Stone space $\mathrm{St}(\mathcal{O}_h)$ of the observable algebra is the
compact Hausdorff space whose points are ultrafilters on $\mathcal{O}_h$, with
the natural Stone embedding $\iota : X \to \mathrm{St}(\mathcal{O}_h)$ sending
$x$ to the principal ultrafilter $\{E \in \mathcal{O}_h : x \in E\}$.

\begin{proposition}[Stone space identification]
\label{II:prop:stone-id}
Suppose $\mathcal{O}_h = \mathcal{B}$ modulo $\mu$.  Then $\iota$ is a
measure-theoretic isomorphism: $\iota$ identifies the measure algebra
$(X/\mu, \mathcal{B}/\mu)$ with
$(\mathrm{St}(\mathcal{O}_h)/\iota_*\mu,
\mathcal{B}(\mathrm{St}(\mathcal{O}_h))/\iota_*\mu)$.
\end{proposition}

\begin{proof}
\emph{Injectivity mod $\mu$.}  Suppose $\iota(x) = \iota(x')$.  For any
$E \in \mathcal{B}$ with $x \in E \not\ni x'$, find $F \in \mathcal{O}_h$
with $\mu(E \triangle F) = 0$.  Since $x \in F$ and $\iota(x) = \iota(x')$,
also $x' \in F$, so $x' \in (E \triangle F) \cup F^c$, placing $x'$ in a
null set.  Hence $\iota$ is injective $\mu$-a.e.

\emph{$\sigma$-algebra identification.}  The pullback $\iota^{-1}$ maps the
Baire $\sigma$-algebra of $\mathrm{St}(\mathcal{O}_h)$ (generated by the
clopen sets $[E] = \{u : E \in u\}$) to $\mathcal{O}_h = \mathcal{B}$ mod
$\mu$.

\emph{Measure recovery.}  Since $\mu$ is $\sigma$-additive on
$\mathcal{O}_h = \mathcal{B}$, the collective exhaustion condition
of~\cite{mahon_paper1} forces $\iota_*\mu$ to concentrate on the principal
ultrafilters $\iota(X)$, so $\iota_*\mu$ is supported on the image and the
measure algebras are identified.
\end{proof}

\noindent The Stone space~\cite{mahon_paper1} — built as a compact ambient
space for deriving $\sigma$-additivity — is, under the reconstruction condition,
the state space $X$ itself.  The construction that began by seeking
a compact completion of the sample space ends by recognising it as the object
the observer was trying to recover.

The classical Takens embedding theorem~\cite{Takens1981} reaches a related
conclusion by a different route: for a $C^2$ diffeomorphism on a compact
$d$-manifold and a generic smooth $h$, the finite-delay map
$\Phi_h^{(k)}(x) = (h(x), h(Tx), \ldots, h(T^{k-1}x))$ is a diffeomorphic
embedding into $\mathbb{R}^{2d+1}$ for $k \geq 2d+1$.  The reconstruction
theorem here requires only measurability, uses all delays simultaneously, and
replaces the dimension count with the algebraic density condition; in exchange,
the embedding is measure-theoretic rather than topological.  On a compact
manifold with generic smooth $h$, the orbit $\{h \circ T^n\}$ separates
points, so Stone--Weierstrass gives $\mathcal{A}_h$ uniformly dense in $C(M)$
and hence $L^2$-dense: the density condition holds and the two results are
compatible.

Reconstruction is an exact asymptotic condition: it asks whether the delay
algebra generates all of $\mathcal{B}$ in the limit of all lags, and says
nothing about rates, finite approximations, or what a finite time series can
certify.  Those questions are the subject of what follows.

% ====================================================================
\ifstandalone
\bibliographystyle{plain}
\bibliography{references}
\fi
